\documentclass{zhvt-classic}

% 使用本地源流明体（GenRyuMin2 TW）按文件路径加载，不依赖系统字体名
% 版框（内框／界栏／书口侧栏／鱼尾）改用与版心印章同一印泥棕色：
%   印章图片 yshs.png 实色 = RGB(135,68,52) = HTML 874434
\usepackage{xcolor}
\usepackage{tikz}
% sealbrown 双注册：xcolor 供 TikZ/pgf，L3 供 cls 版框
\definecolor{sealbrown}{HTML}{874434}
\ExplSyntaxOn
\color_set:nnn { sealbrown } { HTML } { 874434 }
\ExplSyntaxOff

\zhvtset{
  grid color       = sealbrown,
  main font        = GenRyuMin2TW-SB,
  main font path   = C:/Users/Administrator/texmf/fonts/opentype/public/GenRyuMin2TW/,
jiazhu font      = GenRyuMin2TW-M,
  jiazhu font path = C:/Users/Administrator/texmf/fonts/opentype/public/GenRyuMin2TW/,
  % 页码框容纳 4 字：正常页「九十七」+前置空 1 格=4 行恰好；
  % 100+ 页「一百零一」共 5 行超出，末字「一/二/三」被框裁掉（2026-10-04 实测），故放宽到 5
  shukou page number chars = 5,
  % 书口书名框高从默认 6 字减到 4 字，把省出的 2 字高拨给 chapter 框，
  % 使 6 字书名+卷名放进 9 字高的 chapter 框后不再压线（2026-10-05）
  shukou title chars = 4,
}

\title{史记}[孙子吴起列传]
\author{汉·司马迁}
\maker{闲叟漫记·重製}


% ===== 索隐述赞样式（2026-11-08 四次定稿 · 愚大人亲手改定参数）=====
% ⚠ 竖排前提：自定义字体家族必须带 RawFeature={vertical:+vert}；
%   simhei.ttf 加 +vhal 会整体加载失败静默回退，只写 +vert；simfang 同理。
\newCJKfontfamily{\shuzanfont}[Path=C:/Windows/Fonts/,RawFeature={vertical:+vert}]{simhei.ttf}
\newCJKfontfamily{\shuzanbody}[Path=C:/Windows/Fonts/,RawFeature={vertical:+vert}]{simfang.ttf}
% 题头：152x30pt 圆角 6pt 棕底白字；baseline (0,2) 使整框落入界行格子不压版框线；
% 白字节点 x 不补偿、y +8.2pt（像素级墨迹实测 dx/dy≈0，L/R 留白约 20pt）。
\newcommand{\shuzanhead}[1]{%
  \begin{tikzpicture}[baseline={(0,2)}]
    \fill[sealbrown, rounded corners=6pt]
      (-76pt,-15pt) rectangle (76pt,15pt);
    \node[text=white, inner sep=0pt, anchor=center] at (0pt,8.2pt)
      {\shuzanfont\fontsize{23}{23}\selectfont #1};
  \end{tikzpicture}%
}
% ★ 修正书口卷名位置：chapter 框起始位置改回 cls 默认值，避免全局影响。
%   只在每个卷首页通过 \juan 命令里的空命令微调该页位置（安全，不会跨页污染）。
\ExplSyntaxOn
\AtBeginDocument{
  \dim_set:Nn \l__zhvt_shukou_chapter_title_start_dim
    { \l__zhvt_shukou_title_height_dim + \l__zhvt_font_size_dim * 2 }
  \dim_set:Nn \l__zhvt_shukou_page_number_start_dim
    { \l__zhvt_shukou_chapter_title_start_dim + \l__zhvt_shukou_chapter_title_height_dim }
}
\ExplSyntaxOff

\begin{document}

% ★★ 版心（书口）排布覆写：把书名放到鱼尾之下（2026-10-04）
%
%   cls 默认版心自上而下：书名（title 静态框，在鱼尾之上）→ 鱼尾 → 篇名（chapter 动态框）→ 页码
%   本模板改为：          鱼尾（位置不动）→ 书名 → 篇名 → 页码分割线 → 页码
%
%   手法：把 cls 放在 title 静态框里的书名，改放进 chapter 动态框（鱼尾之下）；
%         title 静态框留空。书名取 \l__zhvt_title_str，即 \title{...} 的第一参数，
%         所以换书时只改 \title，不必在此重复写书名。
%
%   ⚠ 三条必须遵守的写法（2026-10-04 实测）：
%     ① 汉字必须用 \hbox{...} 包住。flowfram 会把框内容装进一个宽 14.4mm（= 一个 line_sep
%        = 书口栏宽）的纵向 minipage 并按该宽度折行，而整页靠 PDF /Rotate 90 显示、
%        框内文字本身不变向 —— 不包 \hbox 时「围炉夜话」4 字会被逐字折成 4 行，
%        在版面上就成了横排 4 字并横向溢出版框。
%     ② 竖排下**一个 \hbox 就是一整列**。书名与篇名要同处一列，必须放进同一个 \hbox；
%        分成两个 \hbox 时，后一个会被当成新的一列、甩到隔壁栏去
%        （实测「序」跑到书口栏外侧 14.4mm 的正文列上）。中间的 \quad 在框内是横向空白，
%        显示出来正好是书名与篇名之间的纵向间隔。
%     ③ ★ 内容盒要写成 **与框等宽的定宽盒** `\hbox to \l__zhvt_shukou_chapter_title_height_dim{…\hfil}`。
%        实测：chapter 框是把内容盒沿纵向**居中**摆放的，内容宽度一变，书名就跟着漂：
%        自然宽度的 \hbox（书名 38.4mm）落在 x=157.2mm，加篇名后（57.6mm）落在 x=152.4mm，
%        目录页（无篇名）与正文页的书名就差半格。做成等宽盒后，无论有没有篇名，
%        内容固定从框顶起排，实测两页均落在 x=142.8mm，不再漂。
\ExplSyntaxOn
% 把书口 chapter 框（书名+篇名）起始位置下移 1 字高，避免 4 字卷名顶到鱼尾
% ⚠ cls 的 \__zhvt_calculate_dimensions: 在 \AtBeginDocument 里执行，会重算该 dim；
%   必须把覆写挂到 begin document 之后，否则会被二次覆盖（2026-10-05 实测）。
% (旧方案已移到导言区，直接修改 chapter 框起始位置)
% 卷标题命令 #1=书口篇名(鱼尾) #2=正文卷名（含卷序，卷名内空格用 ~）
%
% ★★ 目录页两项修订（2026-10-05，愚大人指出）★★
%   ① 页码起点统一对齐：
%      原状 —— 目录项写成 \par\hspace{4em}卷名\hspace{4em}页码，页码起点 = 卷名实宽 + 4em。
%      卷名实宽随字数变：「卷一~梁惠王上」6 字 +~、「卷七~离娄上」5 字 +~，
%      差 1 个字位 → 卷七~卷十 4 条的页码整体比其余 10 条高 27.21pt（一个字），看着上下不齐。
%      改法 —— 卷名补全角空格（宽 = 1 字 = 27.215pt）补足到 **7 字符位**再写进目录，
%      14 条页码起点遂全部落在同一条竖线上（起点对齐；末字随页码字数自然参差）。
%      ⚠ 只对目录里的副本 \l_tmpc_tl 补，正文卷名仍用 \l_tmpa_tl 原样 —— 正文像素不受影响。
%   ② 卷名整体上移半字位：前导 \hspace 由 4em 改 3.5em。
%      实测原状目录块在版框内偏下（天头留白 139.6pt / 地脚留白 110.7pt），
%      上移 0.5em = 13.6pt 后上下基本均分，视觉居中。
\NewDocumentCommand { \juan } { mm } {
	\clearpage
	\tl_set:Nn \l__zhvt_chapter_title_tl { #1 }
	\tl_set:Nn \l_tmpa_tl { #2 }
	\tl_set_eq:NN \l_tmpc_tl \l_tmpa_tl
	\int_step_inline:nn { 7 - \str_count:n { #2 } }
	  { \tl_put_right:Nn \l_tmpc_tl { 　 } }
	\tl_set:Nx \l_tmpb_tl { \zhnum { page } }
	\tl_put_left:Nn \l_tmpb_tl { \hspace { 4em } }
	\tl_put_left:Nn \l_tmpb_tl { \l_tmpc_tl }
	\tl_put_left:Nn \l_tmpb_tl { \par \hspace { 3.5em } }
	\addtocontents { toc } { \l_tmpb_tl }
	\skip_horizontal:n { \l__zhvt_font_size_dim * 4 }
	\l_tmpa_tl
	\par
	\thispagestyle { empty }
}
\cs_set_protected:Npn \__zhvt_set_shukou_content: {
	% ★ 鱼尾上方留白（title 静态框）嵌入闲叟钤印：
	%   整页 /Rotate 90，故此处把图用 \rotatebox{90} 转正为竖章；
	%   尺寸卡在版心条幅内（书内约 7mm 宽 × 33.8mm 高，留余量不压鱼尾）。
	\setstaticcontents * { title } {
		\vbox_to_ht:nn { \l__zhvt_baseline_skip_dim } {
			\vskip 14.2pt
			\hbox_to_wd:nn { \l__zhvt_shukou_title_height_dim } {
				\hfil
				\rotatebox [ origin = c ] { 90 } { \includegraphics [ width = 7mm , height = 33.8mm ] { yshs.png } }
				\hfil
			}
		}
	}
% ★ 版心（书口）统一改法（2026-11-16 愚大人定）：18pt＋框内居中＋书口文字同版框色
%   ① 原 27.2pt 时 7 字长篇名（如「廉颇蔺相如列传」）末字溢出 chapter 框、压住页码；
%      改 \fontsize{18pt}{18pt}，前后 \hfil 令字串在框内居中（不随篇名字数漂移）。
%   ② \vbox{\vskip 8.2pt} 定量下移，抵消缩字号后栏宽方向 8.2pt 上偏（像素实测 y 中心 543.3 = 原版）。
%      ⚠ \raisebox / \vspace* 在 flowfram 框内无效；换字号须同步换补偿：16pt→10.1、19pt→7.4。
%   ③ chapter 与 page~number 均置 \color{sealbrown}（HTML 874434），与版框／界行同色；正文仍黑。
%   ⚠ 页码字位、钤印、鱼尾、界行、正文一律未动；五部合订本同此口径。
	\setdynamiccontents * { chapter }
	{ \color{sealbrown} \vbox { \vskip 8.2pt \hbox to \l__zhvt_shukou_chapter_title_height_dim
		{ \hfil \fontsize{18pt}{18pt}\selectfont \l__zhvt_title_str \quad \l__zhvt_chapter_title_tl \hfil } } }
	\setdynamiccontents * { page~number } { \color{sealbrown} \quad \zhnum { page } }
}
\__zhvt_set_shukou_content:
\ExplSyntaxOff

\maketitle

% 单篇，无多卷目录，不渲染目录页

\juan{孙子吴起列传}{卷六十五~孙子吴起列传第五}
\hspace*{2em}孙子武者。齐人也。以兵法见於吴王阖庐。阖庐曰。﹁子之十三篇。吾尽观之矣。可以小试勒兵乎。﹂对曰。﹁可。﹂阖庐曰。﹁可试以妇人乎。﹂曰。﹁可。﹂於是许之。出宫中美女。得百八十人。孙子分为二队。以王之宠姬二人各为队长。皆令持戟。令之曰。﹁汝知而心与左右手背乎。﹂妇人曰。﹁知之。﹂孙子曰。﹁前。则视心。左。视左手。右。视右手。后。即视背。﹂妇人曰。﹁诺。﹂约束既布。乃设鈇钺。即三令五申之。於是鼓之右。妇人大笑。孙子曰。﹁约束不明。申令不熟。将之罪也。﹂复三令五申而鼓之左。妇人复大笑。孙子曰。﹁约束不明。申令不熟。将之罪也。既已明而不如法者。吏士之罪也。﹂乃欲斩左右队长。吴王从台上观。见且斩爱姬。大骇。趣使使下令曰。﹁寡人已知将军能用兵矣。寡人非此二姬。食不甘味。原勿斩也。﹂孙子曰。﹁臣既已受命为将。将在军。君命有所不受。﹂遂斩队长二人以徇。用其次为队长。於是复鼓之。妇人左右前后跪起皆中规矩绳墨。无敢出声。於是孙子使使报王曰。﹁兵既整齐。王可试下观之。唯王所欲用之。虽赴水火犹可也。﹂吴王曰。﹁将军罢休就舍。寡人不原下观。﹂孙子曰。﹁王徒好其言。不能用其实。﹂於是阖庐知孙子能用兵。卒以为将。西破强楚。入郢。北威齐晋。显名诸侯。孙子与有力焉。\\
\hspace*{2em}孙武既死。后百馀岁有孙膑。膑生阿鄄之间。膑亦孙武之后世子孙也。孙膑尝与庞涓俱学兵法。庞涓既事魏。得为惠王将军。而自以为能不及孙膑。乃阴使召孙膑。膑至。庞涓恐其贤於己。疾之。则以法刑断其两足而黥之。欲隐勿见。\\
\hspace*{2em}齐使者如梁。孙膑以刑徒阴见。说齐使。齐使以为奇。窃载与之齐。齐将田忌善而客待之。忌数与齐诸公子驰逐重射。孙子见其马足不甚相远。马有上。中。下。辈。於是孙子谓田忌曰。﹁君弟重射。臣能令君胜。﹂田忌信然之。与王及诸公子逐射千金。及临质。孙子曰。﹁今以君之下驷与彼上驷。取君上驷与彼中驷。取君中驷与彼下驷。﹂既驰三辈毕。而田忌一不胜而再胜。卒得王千金。於是忌进孙子於威王。威王问兵法。遂以为师。\\
\hspace*{2em}其后魏伐赵。赵急。请救於齐。齐威王欲将孙膑。膑辞谢曰。﹁刑馀之人不可。﹂於是乃以田忌为将。而孙子为师。居辎车中。坐为计谋。田忌欲引兵之赵。孙子曰。﹁夫解杂乱纷纠者不控卷。救斗者不搏撠。批亢捣虚。形格势禁。则自为解耳。今梁赵相攻。轻兵锐卒必竭於外。老弱罢於内。君不若引兵疾走大梁。据其街路。冲其方虚。彼必释赵而自救。是我一举解赵之围而收弊於魏也。﹂田忌从之。魏果去邯郸。与齐战於桂陵。大破梁军。\\
\hspace*{2em}后十三岁。魏与赵攻韩。韩告急於齐。齐使田忌将而往。直走大梁。魏将庞涓闻之。去韩而归。齐军既已过而西矣。孙子谓田忌曰。﹁彼三晋之兵素悍勇而轻齐。齐号为怯。善战者因其势而利导之。兵法。百里而趣利者蹶上将。五十里而趣利者军半至。使齐军入魏地为十万灶。明日为五万灶。又明日为三万灶。﹂庞涓行三日。大喜。曰。﹁我固知齐军怯。入吾地三日。士卒亡者过半矣。﹂乃弃其步军。与其轻锐倍日并行逐之。孙子度其行。暮当至马陵。马陵道陕。而旁多阻隘。可伏兵。乃斫大树白而书之曰﹁庞涓死于此树之下﹂。於是令齐军善射者万弩。夹道而伏。期曰﹁暮见火举而俱发﹂。庞涓果夜至斫木下。见白书。乃钻火烛之。读其书未毕。齐军万弩俱发。魏军大乱相失。庞涓自知智穷兵败。乃自刭。曰。﹁遂成竖子之名。﹂齐因乘胜尽破其军。虏魏太子申以归。孙膑以此名显天下。世传其兵法。\\
\hspace*{2em}吴起者。卫人也。好用兵。尝学於曾子。事鲁君。齐人攻鲁。鲁欲将吴起。吴起取齐女为妻。而鲁疑之。吴起於是欲就名。遂杀其妻。以明不与齐也。鲁卒以为将。将而攻齐。大破之。\\
\hspace*{2em}鲁人或恶吴起曰。﹁起之为人。猜忍人也。其少时。家累千金。游仕不遂。遂破其家。乡党笑之。吴起杀其谤己者三十馀人。而东出卫郭门。与其母诀。齧臂而盟曰。﹃起不为卿相。不复入卫。﹄遂事曾子。居顷之。其母死。起终不归。曾子薄之。而与起绝。起乃之鲁。学兵法以事鲁君。鲁君疑之。起杀妻以求将。夫鲁小国。而有战胜之名。则诸侯图鲁矣。且鲁卫兄弟之国也。而君用起。则是弃卫。﹂鲁君疑之。谢吴起。\\
\hspace*{2em}吴起於是闻魏文侯贤。欲事之。文侯问李克曰。﹁吴起何如人哉。﹂李克曰。﹁起贪而好色。然用兵司马穰苴不能过也。﹂於是魏文候以为将。击秦。拔五城。\\
\hspace*{2em}起之为将。与士卒最下者同衣食。卧不设席。行不骑乘。亲裹赢粮。与士卒分劳苦。卒有病疽者。起为吮之。卒母闻而哭之。人曰。﹁子卒也。而将军自吮其疽。何哭为。﹂母曰。﹁非然也。往年吴公吮其父。其父战不旋踵。遂死於敌。吴公今又吮其子。妾不知其死所矣。是以哭之。﹂\\
\hspace*{2em}文侯以吴起善用兵。廉平。尽能得士心。乃以为西河守。以拒秦。韩。\\
\hspace*{2em}魏文侯既卒。起事其子武侯。武侯浮西河而下。中流。顾而谓吴起曰。﹁美哉乎山河之固。此魏国之宝也。﹂起对曰。﹁在德不在险。昔三苗氏左洞庭。右彭蠡。德义不修。禹灭之。夏桀之居。左河济。右泰华。伊阙在其南。羊肠在其北。修政不仁。汤放之。殷纣之国。左孟门。右太行。常山在其北。大河经其南。修政不德。武王杀之。由此观之。在德不在险。若君不修德。舟中之人尽为敌国也。﹂武侯曰。﹁善。﹂\\
\hspace*{2em}吴起为西河守。甚有声名。魏置相。相田文。吴起不悦。谓田文曰。﹁请与子论功。可乎。﹂田文曰。﹁可。﹂起曰。﹁将三军。使士卒乐死。敌国不敢谋。子孰与起。﹂文曰。﹁不如子。﹂起曰。﹁治百官。亲万民。实府库。子孰与起。﹂文曰。﹁不如子。﹂起曰。﹁守西河而秦兵不敢东乡。韩赵宾从。子孰与起。﹂文曰。﹁不如子。﹂起曰。﹁此三者。子皆出吾下。而位加吾上。何也。﹂文曰。﹁主少国疑。大臣未附。百姓不信。方是之时。属之於子乎。属之於我乎。﹂起默然良久。曰。﹁属之子矣。﹂文曰。﹁此乃吾所以居子之上也。﹂吴起乃自知弗如田文。\\
\hspace*{2em}田文既死。公叔为相。尚魏公主。而害吴起。公叔之仆曰。﹁起易去也。﹂公叔曰。﹁柰何。﹂其仆曰。﹁吴起为人节廉而自喜名也。君因先与武侯言曰。﹃夫吴起贤人也。而侯之国小。又与彊秦壤界。臣窃恐起之无留心也。﹄武侯即曰。﹃柰何。﹄君因谓武侯曰。﹃试延以公主。起有留心则必受之。无留心则必辞矣。以此卜之。﹄君因召吴起而与归。即令公主怒而轻君。吴起见公主之贱君也。则必辞。﹂於是吴起见公主之贱魏相。果辞魏武侯。武侯疑之而弗信也。吴起惧得罪。遂去。即之楚。\\
\hspace*{2em}楚悼王素闻起贤。至则相楚。明法审令。捐不急之官。废公族疏远者。以抚养战斗之士。要在彊兵。破驰说之言从横者。於是南平百越。北并陈蔡。卻三晋。西伐秦。诸侯患楚之彊。故楚之贵戚尽欲害吴起。及悼王死。宗室大臣作乱而攻吴起。吴起走之王尸而伏之。击起之徒因射刺吴起。并中悼王。悼王既葬。太子立。乃使令尹尽诛射吴起而并中王尸者。坐射起而夷宗死者七十馀家。\\
\hspace*{2em}太史公曰。世俗所称师旅。皆道孙子十三篇。吴起兵法。世多有。故弗论。论其行事所施设者。语曰。﹁能行之者未必能言。能言之者未必能行。﹂孙子筹策庞涓明矣。然不能蚤救患於被刑。吴起说武侯以形势不如德。然行之於楚。以刻暴少恩亡其躯。悲夫。\\

% ===== 索隐述赞（司马贞《史记索隐》）· 统一样式 =====
% 题头「索隐述赞曰」五字并入反白框；正文每四句一列，定宽盒 \hfil 居列中；
% \fontsize{24}{41} 列距对齐全书界行（121+41k）；\clearpage 使述赞独占一页。
\clearpage
\shuzanhead{索隐述赞曰}\linebreak
{\shuzanbody\fontsize{24}{41}\selectfont
\hbox to \hsize{\hfil 孙子兵法。一十三篇。美人既斩。良将得焉。\hfil}\linebreak
\hbox to \hsize{\hfil 其孙膑脚。筹策庞涓。吴起相魏。西河称贤。\hfil}\linebreak
\hbox to \hsize{\hfil 惨礉事楚。死后留权。\hfil}%
\par}
\end{document}
